\section{Controls: information is the cause}\label{sec:07_controls}

A panel will ask: ``How do you know that the rise in entry comes from information, and not from something
else?'' The paper answers with one theorem and two controls.
\begin{itemize}[leftmargin=2em]
\item The theorem \paper{Prop.~A.3}: if what the challenger learns is held fixed, a stronger incumbent cannot raise
entry.
\item Control 1 \paper{Table~3, Panel~B}: hold the full orders fixed at both strengths. Entry does not rise.
\item Control 2 \paper{Table~3, Panel~C}: hide the price from the challenger. Entry does not rise.
\end{itemize}
In equilibrium, entry rises from the weak economy \(r_0\) to the strong economy \(r_1\). So the rise needs the
change in information that trade produces.

\subsection{Deterrence at a fixed information experiment}\label{sec:07_controls-fixed}

An \emph{information experiment} is the joint law of the class \(j\) and a signal \(S\) that the challenger observes.
With fixed full orders, the price reveals \(S=X\). The theorem allows a general signal. Economy~1 has the stronger
incumbent: \(F_1\le F_0\), a first-order stochastic strengthening.

\begin{resultbox}[Fixed information \paper{Prop.~A.3}]
Fix the joint law of the class \(j\) and the signal \(S\). Fix the class value laws \(G_H,G_L\) and the cost law. In
each economy \(i\in\{0,1\}\), require
\begin{equation}\label{eq:07-product}
\mathcal L_i(\theta,R\mid j,S)=G_j\otimes F_i ,
\end{equation}
and let the cost \(C\) be independent of the class, both values and \(S\). If the incumbent strengthens from \(F_0\)
to \(F_1\), entry does not rise under rule C1 (an indifferent high-cost challenger prepares). Its decrease is exactly
\[
\Ent_0-\Ent_1=\Pr\bigl\{B_{F_1}(\mu(S))<C\le B_{F_0}(\mu(S))\bigr\},\qquad \mu(S)=\Pr(j=H\mid S),
\]
and it is strict exactly when this probability is positive.
\end{resultbox}

Read the product condition \eqref{eq:07-product} as three parts. Each part matters.
\begin{enumerate}[label=(P\arabic*),leftmargin=3em]
\item Given \((j,S)\), the challenger's value \(\theta\) has the class law \(G_j\): the signal tells nothing about
\(\theta\) beyond the class.
\item Given \((j,S)\), the incumbent's value \(R\) has the law \(F_i\): the signal and the class tell nothing about
the incumbent.
\item Given \((j,S)\), the two values are independent.
\end{enumerate}
The independent cost makes \(\mu(S)\) the whole class belief at the moment of decision.

\begin{caution}[Fix the full conditional law]
A fixed class posterior alone does not fix the information about either acquisition value.
Use the product law \eqref{eq:07-product}.
The two examples below show why both conditional marginal laws matter.
\end{caution}

\subsubsection*{The proof, step by step}

\emph{Step 1: the belief.} The challenger decides on \((S,C)\). The cost is independent of \((j,\theta,R,S)\), so
\(\Pr(j=H\mid S,C)=\Pr(j=H\mid S)=\mu(S)\). The joint law of \((j,S)\) is the same in both economies, so
\(\mu(\cdot)\) is the same function in both.

\emph{Step 2: gross profit after the signal.} A prepared challenger keeps \((\theta-\max\{p,R\})_+\); this is zero when
\(\theta<p\) (\Cref{sec:03_auction}). By (P1)--(P3), \((\theta,R)\) has the law \(G_j\otimes F_i\) given
\((j,S)\). So
\[
\E_i\bigl[(\theta-\max\{p,R\})_+\mid j,S\bigr]=\int\!\!\int(y-\max\{p,x\})_+\,dF_i(x)\,dG_j(y)=g_j(F_i),
\]
which does not depend on \(S\). Average over the class with the weights of Step~1:
\[
\E_i\bigl[(\theta-\max\{p,R\})_+\mid S\bigr]=\mu(S)\,g_H(F_i)+[1-\mu(S)]\,g_L(F_i)=B_{F_i}(\mu(S)).
\]

\emph{Step 3: the preparation rule.} Under rule C1 the challenger prepares exactly when \(C\le B_{F_i}(\mu(S))\).

\emph{Step 4: the pointwise indicator difference.} Proposition~1(a) gives \(B_{F_1}(\mu)\le B_{F_0}(\mu)\) at every
belief. Fix one realization, and write \(b_1=B_{F_1}(\mu(S))\le b_0=B_{F_0}(\mu(S))\). If \(C\le b_1\), both
indicators are one. If \(b_1<C\le b_0\), only the first is one. If \(C>b_0\), both are zero. So
\begin{equation*}
\1\{C\le B_{F_0}(\mu(S))\}-\1\{C\le B_{F_1}(\mu(S))\}=\1\{B_{F_1}(\mu(S))<C\le B_{F_0}(\mu(S))\}\ge0 .
\tag*{\paper{Prop.~A.3, proof}}
\end{equation*}

\emph{Step 5: the exact decrease.} The law of \((S,C)\) is the same in both economies: the law of \(S\) is fixed, and
\(C\) is independent of \(S\) with a fixed law. Take expectations of Step~4 under this one law. The left side becomes
\(\Ent_0-\Ent_1\); the right side becomes the probability in the proposition.

\emph{Step 6: a common tie probability.} Let an indifferent high-cost challenger prepare with a probability \(q\) that
is the same in both economies. Each indicator becomes \(\1\{C<B\}+q\,\1\{C=B\}\). With \(b_1\le b_0\), the difference
is \(0\) for \(C<b_1\) and for \(C>b_0\); \(1-q\) at \(C=b_1<b_0\); \(1\) for \(b_1<C<b_0\); \(q\) at \(C=b_0>b_1\); and
\(0\) when \(b_0=b_1\). Each value is nonnegative, so entry does not rise here either.

\begin{intuition}[Why fixed information gives deterrence]
Hold fixed what the challenger can learn. Then its belief after each signal is the same in both economies. A stronger
rival lowers gross profit at that belief. It can only remove a reason to prepare; it never adds one. So any rise in
entry must come through a change in what the challenger learns.
\end{intuition}

\subsubsection*{Why each part of the condition matters: two examples}

Each example removes one restriction, (P1) or (P2). In both, the joint law of the class and the
signal is fixed, the class value laws are fixed, the cost is independent and the incumbent strengthens. Entry still
rises, because the signal carries information about a value. \Cref{sec:benchmark-controls} gives the numbers of
both examples.

\paragraph{Example A: the signal tells the challenger about the incumbent (breaks P2).} The incumbent value is a
half-half mixture of a low component and a high component. The stronger incumbent moves the high component up, so
\(F_1\le F_0\). Both point-mass classes lie above every incumbent value, so the challenger always wins and pays
\(\max\{p,R\}\). In economy~0 a fair binary signal reports the component of \(R\). In economy~1 the signal is
independent of \(R\). In both, the signal is independent of the class, so \(\mu(S)\) stays at the prior.
\begin{itemize}[leftmargin=2em]
\item Economy~0: after the low component, gross profit is high and a high-cost challenger prepares. After the high
component, gross profit falls below \(c_H\), and it stays out.
\item Economy~1: the signal does not reveal the high component. The average gross profit clears \(c_H\) after every
signal, so a high-cost challenger always prepares.
\end{itemize}
Unconditional gross profit falls, as Proposition~1 requires. But entry rises: the signal in economy~0 told the
challenger when the rival was strong, and economy~1 removes that warning.

\paragraph{Example B: the signal tells the challenger about its value within a class (breaks P1).} Each class
value law has two point masses, a lower value and an upper value. The class \(j\) and the signal \(S\) are
independent fair binary variables, so \(\mu(S)\) stays at the prior. The incumbent is uniform and independent of every
other draw, so (P2) and (P3) hold. Every value exceeds every incumbent value.
\begin{itemize}[leftmargin=2em]
\item Economy~0: \(\theta\) is independent of \(S\) within each class. Gross profit after each signal lies just below
\(c_H\), so only a zero-cost challenger prepares.
\item Economy~1: the incumbent is a little stronger, and \(S\) selects the lower or the upper value of the class. After
the upper signal, gross profit clears \(c_H\), and a high-cost challenger prepares.
\end{itemize}
The class gross profits fall, but entry rises. The class belief never moves; the value information does the work.

\paragraph{The third part and the cost.} Part (P3) lets Step~2 integrate \(R\) against \(F_i\) at each \(\theta\). If
the values were correlated given the signal, the rent after the signal would not be \(g_j(F_i)\). A cost that is
informative about the class or the values would act as a second signal; the condition would then have to hold for
\((S,C)\) in place of \(S\).

\subsection{Control 1: full orders held fixed}\label{sec:07_controls-frozen}

Proposition~A.3 compares economies with the same experiment. On the equilibrium path the experiment changes with
\(r\) (no trade at \(r_0\), full orders at \(r_1\)), so the theorem says nothing about that comparison. The two
controls build the comparison that the theorem covers.

The control fixes \(\lambda\), \(\beta\), \(\pi=1/2\), the cost law and the full orders \((\delta_1,\delta_{-1})\) at
every strength, and lets the challenger and the market makers respond \paper{App.~A.6}.

\emph{Step 1: posteriors.} Under full orders, \(f_H=(b,b_0,b+\lambda)\) and \(f_L=(b+\lambda,b_0,b)\) on \((-1,0,1)\),
with \(b=(1-\lambda)\beta\) and \(b_0=(1-\lambda)(1-2\beta)\). Bayes' rule gives \((m,\tfrac12,M)\)
(\Cref{sec:04_slot_beliefs}).

\emph{Step 2: the price reveals the order.} Entry is \(e_r(\mu)=\rho+(1-\rho)\1\{B_r(\mu)\ge c_H\}\) and the price is
\(P(x)=t_0+e_r(\mu(x))\,w(\mu(x))\), with \(w(\mu)=t_L-t_0+\DT\mu\). Entry is nondecreasing in the posterior, so the
three prices differ (\Cref{sec:04_slot_beliefs}). So the signal is \(S=X\), and its law given the class, \(f_j\),
does not depend on \(r\).

\emph{Step 3: the product condition holds.} The order depends only on the class and the slot draws, and the slot
draws are independent of \((\theta,R)\). Given \(j\), \(\theta\sim G_j\) and \(R\sim F_r\) are independent. So
\(\mathcal L(\theta,R\mid j,X)=G_j\otimes F_r\), and Proposition~A.3 applies to every pair of strengths.

\emph{Step 4: which posterior recruits.}
\Cref{sec:03_auction-derivatives} derives \(\tau\) and proves that it rises. It lies above the prior already at the
start of the benchmark path. Thus neither a sale nor zero recruits on the benchmark path.
A purchase recruits exactly when \(\tau(r)\le M\), or \(r\le r_C\), the entry ceiling.
The equality uses rule C1; \Cref{sec:08_coexistence-ties} gives the alternative at \(r_C\).

\emph{Step 5: entry.} A purchase has probability \(\tfrac12(b+\lambda)+\tfrac12b=b+\lambda/2\). So frozen
entry is \(\rho+(1-\rho)(b+\lambda/2)\) on the benchmark path up to \(r_C\), and \(\rho\) for \(r>r_C\).

\emph{Step 6: the crossing probabilities.} Check \(\{B_{r_b}(\mu(X))<C\le B_{r_a}(\mu(X))\}\) order by order.
\begin{itemize}[leftmargin=2em]
\item \(r_0\to r_1\). A zero cost prepares in both economies. After a purchase, a high cost prepares in both, because
\(B_{r_1}(M)\ge c_H\) and \(B_{r_0}(M)\ge B_{r_1}(M)\). After a sale or the zero order it prepares in neither, because
\(B_{r_0}(1/2)<c_H\). The event is empty: its probability is zero.
\item \(r_1\to r_2\). The event is \(\{X=1,\ C=c_H\}\), because \(B_{r_2}(M)<c_H\le
B_{r_1}(M)\). Its probability is \((1-\rho)(b+\lambda/2)\), exactly the fall in entry.
\end{itemize}

\emph{Step 7: at \(r_0\) the frozen profile is a control, not an equilibrium.} There \(e(1)=1\), because
\(B_{r_0}(M)\ge B_{r_1}(M)\ge c_H\). The informed orders earn
\begin{align*}
\text{purchase: }&\DT(r_0)(1-M)-k,\\
\text{sale: }&\rho\,\DT(r_0)\,m-k.
\end{align*}
Each gross return is below \(\DT(r_0)\), and (H2) gives \(\DT(r_0)<k\). Both net returns are negative, so the investor
would abstain. At \(r_1\) the frozen tuple is the unique equilibrium, because the
trading bound forces full orders there (\Cref{sec:06_main_result}).

\begin{panelnote}[Flat entry under fixed information]
The frozen experiment gives only three posteriors.
Only the top posterior can recruit on this path.
A higher threshold leaves entry unchanged until it exceeds that posterior.
Thus frozen entry stays at one level through \(r_C\), then falls to the floor \(\rho\).
The old Laplace experiment had a continuum of intermediate posteriors.
Its frozen entry could fall as the threshold moved within that continuum.
\end{panelnote}

\subsection{Control 2: the hidden price}\label{sec:07_controls-hidden}

Now the challenger cannot see the price; the investor and the market makers act as before.
\begin{enumerate}[leftmargin=2em]
\item \emph{Entry.} The challenger's belief stays at \(1/2\). The high-cost window (H1) gives \(B_{r_0}(1/2)<c_H\), and
Proposition~1(a) gives \(B_{r_1}(1/2)\le B_{r_0}(1/2)\). So entry is \(\rho\) after every order at both strengths.
\item \emph{Orders at \(r_0\).} The gross return of a correctly signed nonzero order is at most \(\DT(r_0)\), and (H2)
gives \(\DT(r_0)<k\). A wrong-signed order earns at most \(-k\). Both classes abstain.
\item \emph{Orders at \(r_1\).} The trading bound holds for any entry in \([\rho,1]\), so with \(e\equiv\rho\). The gross
returns of the two correctly signed nonzero orders satisfy \(K_H=\rho\DT(r_1)[1-\mu(1)]\ge\rho\DT(r_1)m\) and
\(K_L=\rho\DT(r_1)\mu(-1)\ge\rho\DT(r_1)m\). At the prior one half, \(m=\kappa(1/2)\), so (H2) gives
\(\rho m\DT(r_1)>k\). Both classes trade fully. The sale margin equals that of the equilibrium at \(r_1\), because a
sale meets entry \(\rho\) in both games.
\item \emph{Outcome.} At \(r_1\) the three prices differ: the price is informative, but the challenger does not see
it. Entry is \(\rho\), high-class ownership is \(\rho/2\), and proceeds are \(t_0+\rho\,w(1/2)\) at both strengths.
\end{enumerate}
Each tuple is the unique equilibrium of the game in which the challenger decides on its prior.
\Cref{sec:benchmark-controls} gives the benchmark values of both controls.

\begin{resultbox}[The identification logic]
At fixed information, entry cannot rise with incumbent strength.
The equilibrium rise therefore needs the change in information that trade produces.
The two controls isolate this change: frozen full orders keep entry flat; a hidden price keeps entry at its floor.
\end{resultbox}
